\answer{ex:23:1}{Ответ падает как $(1-f)^n$; вдвое --- после $n=\ln2/[-\ln(1-f)]\approx14$, $7$ и $1$ запаха (модель при $f=5\%$: $0.50$ после $14$). Без забывания банк быстро «забивается»; $\tau_c$ освобождает строки.}
\answer{ex:23:2}{$K=5$: $\binom{2000}{2}/\binom{51}{5}\approx0.85$ пары (при $K=4$ --- $8$, при $K=7$ --- $0.017$). Доля общих входов $K/51$: $0.14$ при $K=7$ и $0.49$ при $K=25$ --- клетки становятся полукопиями друг друга.}
\answer{ex:23:3}{$\ln2/[-\ln(1-\eta)]$: $6.6$ постукивания ($66$~с) без еды и $13.5$ ($135$~с) на еде. Конечный уровень $A_\infty$ тот же: флаг меняет скорость, а не цель.}
\answer{ex:23:4}{$\Delta w/(\eta\tau)=e^{\Delta/\tau}(1-e^{-T/\tau})$ при $\Delta<0$; $e^{-\Delta/\tau}-e^{-(T-\Delta)/\tau}$ при $0\le\Delta\le T$; $-(1-e^{-T/\tau})\,e^{-(\Delta-T)/\tau}$ при $\Delta>T$. Ноль --- при $\Delta=T/2$.}
\answer{ex:23:5}{Одно сочетание: $0.95$, $0.17$, $0.034$; десять: $1.40$, $0.50$, $0.27$. На четвёртый день память почти целиком в медленном банке, и десять сочетаний дают её в восемь раз больше, чем одно.}
\answer{ex:23:6}{С неподвижным порогом загорается $\approx65\%$ клеток, и выученный запах избегается лишь на $\approx0.08$; при $1.5$-кратной силе --- $36\%$ и $0.14$. С отбором $k$ победителей код не меняется, отвращение $1.00$.}
\answer{ex:23:7}{Открытая задача. Выход, выученный на сахар, должен возбуждать учителя наказания, а сахар --- гасить это возбуждение; тогда учитель наказания несёт отрицательную часть $\delta$ из~\eqref{eq:ctd}, учитель награды --- положительную, и пара линий заменяет одно число со знаком.}
